\section{Introduction}
% C101: integration-level distinction; E101 abstract only. Remaining wording defines this study's design scope.
Cyber-physical production systems connect material-handling equipment, sensing, and control software to coordinate physical operations with computational decisions. In the cell considered here, production continuity depends on knowing whether a workpiece has reached a transfer point, whether a pneumatic movement has completed, and whether the robot can execute the next handling operation. A digital representation becomes useful for this purpose when its state estimates and predictions can support those decisions. Kritzinger \emph{et al.} classify manufacturing representations as digital models, digital shadows, and digital twins according to their level of integration~\cite{Kritzinger2018}. This distinction motivates the scope adopted in this study: an executable physical model forms the initial simulation basis, while observation processing, synchronization, prediction, and supervisory action are proposed as additional layers.

% C102: local project scope and fault examples; E107 and E108. No experimental outcome is asserted.
The target application is a sorting and pick-and-place cell comprising two conveyors, a pneumatic pusher, a Cartesian manipulator, and a vacuum gripper. The intended sensing arrangement includes photoelectric workpiece detection and end-position feedback for pneumatic motion. A sequential controller advances through feeding, transfer, tracking, pickup, placement, and return states. For this application, a missing detection event is ambiguous: the workpiece may not have arrived, the sensor may have failed, or an observation may have been delayed or lost in transmission. A supervisory response therefore needs to account for both equipment behavior and the validity of the information used to infer that behavior. Continuing a sequence on an unsupported state estimate and stopping whenever an observation is missing represent different operating choices whose consequences must be evaluated under the same disturbances.

% C103: networked estimation and information timeliness; E102 Sec. I(b)-(c), E103 abstract and Sec. I.
Communication effects are part of the control problem when state information crosses a network. Hespanha \emph{et al.} review estimation and control under limitations that include sampling, variable delay, and packet dropouts~\cite{Hespanha2007}. Yates \emph{et al.} examine the age of information as a measure of the timeliness of timestamped status updates~\cite{Yates2021}. These perspectives motivate treating the age and delivery history of an observation as inputs to the proposed twin. For the present cell, a workpiece estimate derived from an old photoelectric event should carry a different confidence level from one corrected by a recent event. The implementation will therefore distinguish source timestamps, reception times, sequence numbers, and clock domains. Uncertainty in clock alignment will be retained when interpreting observation age, rather than assuming that timestamps from different devices share an exact time base.

% C104: published abstract supports diagnosis/inverse-model architecture and simulated material-handling case; E104 abstract-only.
Sensor-fault accommodation has also been investigated for automated manufacturing. El-Mahdy \emph{et al.} describe an active fault-tolerant approach using separate long short-term memory networks for diagnosis and inverse-model-based control reconfiguration, evaluated through a Factory I/O--MATLAB material-handling simulation~\cite{ElMahdy2024}. Their study provides a relevant example of linking diagnosis to a revised control action. In the present project, the initial design will use interpretable residuals, event-consistency checks, and explicit uncertainty limits. This choice allows a missed workpiece event, an overdue cylinder transition, or stale data to be connected to a stated operating assumption before considering more complex learned models.

% C105: published author-hosted journal version; E105 p. 1 abstract and Sec. 1.
Predictive digital-twin control provides a complementary basis for evaluating future actions. Chen \emph{et al.} combine multistep prediction and quantile-based uncertainty estimation with robust model predictive control, including an additive-manufacturing case study~\cite{Chen2026}. That work motivates the joint consideration of prediction uncertainty and computation time. The handling problem examined here instead concerns the timing of discrete production events and selection among a finite set of supervisory policies. Its design question is whether prediction remains useful when the observations initializing it may be faulty or stale and the selected response must be returned before the next handling decision.

% C106: scoped research question and proposed design, not a literature-wide absence or novelty claim; E108.
The research question is consequently framed as follows: under which sensing and communication conditions can a network-aware predictive twin preserve task completion by selecting an admissible software policy before its decision deadline? Answering this question requires coupling state confidence, fault hypotheses, prediction cost, and the available control actions. A geometrically detailed simulation alone does not answer it. The specific integration problem addressed in this project is to carry the uncertainty attached to received observations through diagnosis and short-horizon scenario evaluation to a supervisory decision. The study will also examine when the remaining observations are insufficient to support continued handling and a controlled hold is the appropriate response.

% C107: MuJoCo capabilities, E106 abstract; implementation boundary, E107; proposed architecture, E108.
To investigate this question, the proposed framework builds on the existing MuJoCo cell model. MuJoCo provides multibody and contact-dynamics functionality intended for model-based control~\cite{Todorov2012}. Within the planned experimental architecture, one simulator instance represents the plant, while a separate predictive instance is initialized from sensor-derived estimates. Simulator ground truth is reserved for evaluation. A hybrid estimator will combine continuous motion estimates with discrete task states, and a fault-monitoring layer will assess residuals together with event order, timeouts, and data freshness. Candidate rollouts will explore plausible disturbances and permitted responses within a fixed computation budget. The supervisor will issue a response only while its state assumptions and command validity interval remain applicable. Local sequence control and interlocks will retain authority over actuator execution.

% C108: proposed contributions; E108. These are commitments for future implementation, not completed findings.
The proposed contributions are threefold:
\begin{enumerate}
  \item An architecture linking timestamped observations, hybrid state estimates, fault hypotheses, and observation confidence to supervisory decisions for a conveyor-based handling cell.
  \item A prediction-and-selection mechanism that evaluates admissible operating policies under sensing and communication disturbances while enforcing a computation deadline and retaining a defined hold response.
  \item An evaluation protocol comparing sequential control, fault-aware supervision, and network-aware predictive supervision using matched initial conditions and disturbance schedules, with separate measures of estimation quality, detection performance, production completion, recovery, and computational timeliness.
\end{enumerate}

% C109: project capability and validation limits; E107 and E108.
In this study, self-reconfiguration refers to changes in software policies supported by the existing equipment: revising conveyor timing, restricting the admission of new workpieces, selecting valid observations, or holding the sequence. Continuous speed adjustment is an admissible response only when the installed conveyor drive supports it. The current prototype supplies the physical simulation and a commissioning sequence; the estimator, communication model, predictive supervisor, and their integrated validation remain under development. Calibration against the physical cell is also pending. Accordingly, this research-in-progress draft presents the problem formulation and intended methodological contributions without claiming measured improvements in throughput, fault-detection accuracy, or recovery performance.
